\section*{Capítulo \ref{ch:b1:geological-maps} --- Mapas geológicos}

\begin{solution}{exo:b1:geological-maps:1}
\omterm{def:b1:geological-maps:attitude}{Rumbo}: \omterm{def:b1:spherical-astronomy:horizontal}{acimut} de las rectas horizontales del plano; \omterm{def:b1:geological-maps:attitude}{buzamiento}: ángulo entre el plano y la
horizontal, medido perpendicularmente al \omterm{def:b1:geological-maps:attitude}{rumbo}; \omterm{def:b1:geological-maps:attitude}{dirección de buzamiento}: \omterm{def:b1:spherical-astronomy:horizontal}{acimut} de la recta de
máxima pendiente descendente. Para $135/40$, el \omterm{def:b1:geological-maps:attitude}{rumbo} según la regla de la mano derecha es
$135 - 90 = \ang{045}$. Símbolo: una línea de \omterm{def:b1:geological-maps:attitude}{rumbo} orientada N\,45\,E, con un trazo en su lado
sureste (hacia \ang{135}) etiquetado con 40.
\end{solution}

\begin{solution}{exo:b1:geological-maps:2}
$z = 500 - \tan\ang{30}\,(x\sin\ang{90} + y\cos\ang{90}) = 500 - 0.577\,x$.
En $(1000, 0)$: $500 - 577 = \qty{-77}{m}$. En $(0, 1000)$: \qty{500}{m}, ya que el plano es
horizontal a lo largo de su \omterm{def:b1:geological-maps:attitude}{rumbo} norte--sur.
\end{solution}

\begin{solution}{exo:b1:geological-maps:3}
Una traza que sigue una curva de nivel es un contacto horizontal, a \qty{350}{m}. Una traza que
atraviesa en línea recta colinas y valles pertenece a un contacto vertical.
\end{solution}

\begin{solution}{exo:b1:geological-maps:4}
A lo largo del río, el núcleo del \omterm{def:b1:geological-maps:fold}{anticlinal} muestra el esquisto y el granito (bajo la arenisca);
el núcleo del \omterm{def:b1:geological-maps:fold}{sinclinal} lo ocupa la creta. Las unidades más antiguas quedan más abajo en la
sucesión: donde las capas se arquean hacia arriba y la erosión corta más hondo, aparecen las más
antiguas; donde se hunden, se conservan las más jóvenes en el centro.
\end{solution}

\begin{solution}{exo:b1:geological-maps:5}
Los dos primeros puntos están a la misma altitud: la recta $y = 0$ es una línea de \omterm{def:b1:geological-maps:attitude}{rumbo}, así que
el \omterm{def:b1:geological-maps:attitude}{rumbo} es este--oeste. El contacto está \qty{100}{m} más bajo \qty{500}{m} al norte: buza hacia el
norte, $\tan\delta = 100/500 = 0.2$, $\delta \approx \ang{11.3}$; \omterm{def:b1:geological-maps:attitude}{dirección de buzamiento}
\ang{000}, \omterm{def:b1:geological-maps:attitude}{rumbo} \ang{270} según la regla de la mano derecha. Con el teorema:
$\overrightarrow{P_1P_2} = (1000, 0, 0)$,
$\overrightarrow{P_1P_3} = (0, 500, -100)$, $\vec n = (0, 100\,000,
500\,000)$, $\cos\delta = 500\,000/509\,902 \approx 0.981$, $\delta \approx
\ang{11.3}$, y $(n_x, n_y) = (0, 100\,000)$ apunta al norte: el mismo resultado.
\end{solution}

\begin{solution}{exo:b1:geological-maps:6}
Aguas abajo es el norte, $s = 0.03$ (\ang{1.72}). (a) $a = 0 < s$: la V apunta aguas arriba
(al sur), a lo largo de las curvas de nivel. (b) $a = -\tan\ang{20} < 0 < s$: V aguas arriba.
(c) $a = \tan\ang{1} \approx 0.017 < s$: V aguas arriba. (d) $a = \tan\ang{20}
\approx 0.364 > s$: V aguas abajo (al norte).
\end{solution}

\begin{solution}{exo:b1:geological-maps:7}
$\xi_{\text{fondo}} - \xi_{\text{flanco}} = h/(a - s) = 80/(0.364 - 0.03)
\approx \qty{240}{m}$: la traza en el flanco queda unos \qty{240}{m} aguas arriba del vértice.
\end{solution}

\begin{solution}{exo:b1:geological-maps:8}
El cierre del \omterm{def:b1:geological-maps:fold}{anticlinal} apunta en la dirección de la inmersión: el eje se hunde hacia el norte.
Con un eje que se hunde hacia el norte, los \omterm{def:b1:geological-maps:geological-map}{afloramientos} de un \omterm{def:b1:geological-maps:fold}{sinclinal} se cierran hacia el
sur.
\end{solution}

\begin{solution}{exo:b1:geological-maps:9}
El lado sureste-este-sureste muestra unidades más jóvenes (creta y lutita frente a unidades más
antiguas situadas enfrente): ese bloque se hundió. Los trazos del símbolo de la falla están en ese
lado, el bloque hundido.
\end{solution}

\begin{solution}{exo:b1:geological-maps:10}
Una recta común a los dos planos es perpendicular a las dos normales, luego paralela a
$\vec n_1 \times \vec n_2$. Para $090/30$: $\vec n_1 = (0.577, 0,
1)$; para $270/30$: $\vec n_2 = (-0.577, 0, 1)$. Entonces $\vec n_1 \times \vec n_2
= (0\cdot 1 - 1\cdot 0,\ 1\cdot(-0.577) - 0.577\cdot 1,\ 0) = (0, -1.155,
0)$: una recta norte--sur horizontal. El primer plano buza hacia el este, el segundo hacia el oeste:
buzan alejándose de la recta, que es la cresta de un \omterm{def:b1:geological-maps:fold}{anticlinal} de eje horizontal.
\end{solution}

\begin{solution}{exo:b1:geological-maps:11}
La traza de un dique vertical es la recta que pasa por el origen con \omterm{def:b1:geological-maps:attitude}{rumbo} \ang{045}, sea cual sea
la altitud: en la colina queda sobre esa misma recta. El dique buzante es el plano $z = -\tan\ang{60}\,d$,
donde $d$ es la distancia horizontal a la recta de \omterm{def:b1:geological-maps:attitude}{rumbo} que pasa por el origen, contada positiva
hacia el sureste. A $z = \qty{300}{m}$: $d = -300/\tan\ang{60}
\approx \qty{-173}{m}$: la traza queda a \qty{173}{m} al noroeste de la recta de \omterm{def:b1:geological-maps:attitude}{rumbo} que pasa por
el origen, del lado de la pendiente hacia arriba.
\end{solution}

\begin{solution}{exo:b1:geological-maps:12}
\[
  \frac{\partial z}{\partial x} \approx \frac{405 - 420}{100} = -0.15, \qquad
  \frac{\partial z}{\partial y} \approx \frac{412 - 420}{100} = -0.08,
\]
así que la pendiente es $\arctan\sqrt{0.15^2 + 0.08^2} = \arctan 0.17 \approx
\ang{9.6}$. El terreno desciende a lo largo de $-\nabla z = (0.15, 0.08)$, cuyo \omterm{def:b1:spherical-astronomy:horizontal}{acimut} es
$\arctan(0.15/0.08) \approx \ang{62}$: la celda mira al este-nordeste.
\end{solution}

\begin{solution}{pb:b1:geological-maps:1}
\textbf{1.} Esquisto, granito, arenisca, caliza, lutita, creta, aluvión.

\textbf{2.} De oeste a este: creta y lutita, caliza, arenisca, luego esquisto y granito cerca de
$x = \qty{3.5}{km}$, arenisca, caliza y lutita otra vez, después creta; aluvión a lo largo de todo
el fondo.

\textbf{3.} La arenisca, con el esquisto y el granito bajo ella a lo largo del río: las unidades más
antiguas en el núcleo, un \omterm{def:b1:geological-maps:fold}{anticlinal}.

\textbf{4.} La creta, la unidad más joven: un \omterm{def:b1:geological-maps:fold}{sinclinal}.

\textbf{5.} La base de la arenisca es más alta sobre el eje del \omterm{def:b1:geological-maps:fold}{anticlinal}, y solo a lo largo del río
la erosión ha cortado por debajo de ella.

\textbf{6.} Una discordancia: la arenisca descansa sobre dos unidades antiguas distintas, esquisto y
granito, y corta el contorno del granito.

\textbf{7.} $435 - 261 = 174$, $432 - 428 = 4$, $512 - 586 = \qty{-74}{m}$.

\textbf{8.} $\overrightarrow{E_1E_2} = (600, 1200, -170)$,
$\overrightarrow{E_1E_3} = (1100, 400, -248)$.

\textbf{9.} El producto vectorial es
\[
  (-229\,600,\ -38\,200,\ -1\,080\,000);
\]
orientado hacia arriba, $\vec n_E = (229\,600,\ 38\,200,\ 1\,080\,000)$.

\textbf{10.} $\|\vec n_E\| \approx 1\,104\,800$; $\cos\delta \approx 0.9776$,
$\delta \approx \ang{12.2}$.

\textbf{11.} \omterm{def:b1:spherical-astronomy:horizontal}{Acimut} de $(229\,600, 38\,200)$: $\arctan(229\,600/38\,200)
\approx \ang{80.6}$; \omterm{def:b1:geological-maps:attitude}{actitud} $081/12$, \omterm{def:b1:geological-maps:attitude}{rumbo} $\approx \ang{351}$.

\textbf{12.} Sí: el símbolo del flanco oriental tiene una línea de \omterm{def:b1:geological-maps:attitude}{rumbo} casi norte--sur y un trazo
de \omterm{def:b1:geological-maps:attitude}{buzamiento} hacia el este, etiquetado con 12.

\textbf{13.} $524 - 669 = -145$, $421 - 459 = -38$, $319 - 209 =
\qty{110}{m}$.

\textbf{14.} $\overrightarrow{W_1W_2} = (700, 1200, 107)$ y
$\overrightarrow{W_1W_3} = (1100, -600, 255)$. Su producto vectorial, orientado hacia arriba, es
\[
  (370\,200,\ -60\,800,\ -1\,740\,000) \quad\longrightarrow\quad
  \vec n_W = (-370\,200,\ 60\,800,\ 1\,740\,000).
\]

\textbf{15.} Así, $\cos\delta = 1\,740\,000/\|\vec n_W\| \approx 0.9776$, $\delta
\approx \ang{12.2}$; \omterm{def:b1:geological-maps:attitude}{dirección de buzamiento}: \omterm{def:b1:spherical-astronomy:horizontal}{acimut} de $(-370\,200, 60\,800)$,
$\approx \ang{279.3}$. \omterm{def:b1:geological-maps:attitude}{Actitud} $279/12$, \omterm{def:b1:geological-maps:attitude}{rumbo} $\approx \ang{189}$: el flanco occidental buza hacia
el oestenoroeste, y el flanco oriental, hacia el estenordeste.

\textbf{16.} $\vec n_E \times \vec n_W = (804 \times 10^6,\ -799\,320 \times
10^6,\ 28\,101 \times 10^6)$ (redondeado).

\textbf{17.} La dirección hacia abajo es $(-804, 799\,320, -28\,101) \times
10^6$: \omterm{def:b1:geological-maps:attitude}{rumbo} $\arctan(-804/799\,320) \approx \ang{0}$, es decir, el norte; inmersión
$\arctan\bigl(28\,101/\sqrt{804^2 + 799\,320^2}\bigr) \approx \ang{2.0}$. El eje se hunde
\ang{2} hacia el norte.

\textbf{18.} $8000 \times \tan\ang{2.0} \approx \qty{280}{m}$ más bajo, hacia
\qty{-50}{m}. El núcleo del \omterm{def:b1:geological-maps:fold}{anticlinal} es más alto cerca del río y se hunde hacia el norte, así que
su \omterm{def:b1:geological-maps:geological-map}{afloramiento} se cierra hacia el norte, como muestra el mapa.

\textbf{19.} $\overrightarrow{F_1F_2} = (1627, 2500, -276)$ y
$\overrightarrow{F_1F_3} = (3124, 5500, 77)$ dan la normal hacia arriba
\[
  (1\,710\,500,\ -987\,503,\ 1\,138\,500), \qquad
  \delta = \arccos\frac{1\,138\,500}{2\,279\,700} \approx \ang{60},
\]
y una \omterm{def:b1:geological-maps:attitude}{dirección de buzamiento} $\approx \ang{120}$. \omterm{def:b1:geological-maps:attitude}{Actitud} $120/60$, \omterm{def:b1:geological-maps:attitude}{rumbo} \ang{030}.

\textbf{20.} La falla buza hacia el estesureste, así que ese lado es el bloque superior, el que queda
encima del plano de falla buzante; se hundió (unidades más jóvenes en superficie): es una falla normal.

\textbf{21.} El plano de falla buza hacia el estesureste, aguas abajo y mucho más deprisa que
desciende el fondo del río: el terreno más bajo corta el plano más abajo en la pendiente, así que la
traza gira hacia el estesureste en el valle --- la regla de la V aplicada al plano de falla.

\textbf{22.} Intrusión del granito; erosión antes de la arenisca; sedimentación de la creta;
plegamiento; fallamiento; sedimentación del aluvión.

\textbf{23.} Después de hace 90 millones de años (la creta está plegada) y antes de hace 2.58
millones de años (el aluvión no está fallado); el plegamiento vino primero, y el fallamiento después.

\textbf{24.} La falla desciende \qty{1405}{m} desde F1 (\qty{405}{m}) hasta \qty{-1000}{m}; con un
\omterm{def:b1:geological-maps:attitude}{buzamiento} de \ang{60}, eso supone una distancia horizontal, en la pendiente, de
$1405/\tan\ang{60} \approx \qty{811}{m}$. El pozo debe perforarse unos
\textbf{\qty[text-series-to-math=true]{810}{m} al estesureste de F1} (\omterm{def:b1:spherical-astronomy:horizontal}{acimut} \ang{120}).
\end{solution}
